\documentclass[12pt]{article}
\usepackage[utf8]{inputenc}
\usepackage[T1]{fontenc}
\usepackage{hyperref}
\title{Wongo Language Resources for K-12 Curriculum Development}
\author{Compiled by Claude Sonnet 5.5 (claude-sonnet-5-5) for Andr\'as Kornai}
\date{}
\begin{document}
\maketitle

\section{Introduction}
Wongo (ISO 639-3 \texttt{won}) is a Bantu language of Kasai, Democratic Republic of the Congo (Guthrie C.85, in the Bushoong C.80 group). It must not be confused with Bushoong (\texttt{buf}), the language of the Kuba Kingdom, which is a separate language. Wikipedia gives about 13,000 speakers (2000 figure), and a University of Missouri profile gives about 17,000. Joshua Project reports about 29,000 in its Wongo people-group profile. Mostly Roman Catholic, the Wongo are culturally influenced by the Kuba and Lele. Very little material exists.

\section{Primary Educational Resources}
No primers, textbooks or school programmes in Wongo were found.

\section{Linguistic and Reference Materials}
\begin{itemize}
\item Wikipedia, Wongo language (a stub): \url{https://en.wikipedia.org/wiki/Wongo_language}
\item University of Missouri DICE profile of Wongo: \url{https://dice.missouri.edu/assets/docs/niger-congo/Wongo.pdf}
\end{itemize}

\section{Religious and Community Resources}
\begin{itemize}
\item Joshua Project, Wongo in Congo, Democratic Republic of: \url{https://joshuaproject.net/people_groups/15416}
\item Scripture Earth page for Wongo, which lists a Gospel of John (Bible.com entry) in Kiwongo. I did not see details. \url{https://scriptureearth.org/00eng.php?iso=won}
\end{itemize}

\section{Digital and Audio Resources}
Scripture Earth states that Global Recordings Network audio exists for Wongo, but I saw no direct URL for it. The Global Recordings page seen in search results is for Bushongo (\texttt{buf}), a different language, and is not cited as a Wongo resource.

\section{Curriculum Development Context}
Wongo is not known to be taught in schools. Documentation is limited, with no orthography or school material found. The sources are dated or secondary. A curriculum would need community-led orthography work, and French and the regional languages would remain the languages of wider communication.
\end{document}